problem:  
Let \(M^{2n}\) be a closed, simply connected Riemannian manifold with nonnegative sectional curvature that admits an isometric and effective torus action \(T^k \curvearrowright M\). Assume \(k \ge \frac{n}{2}\) and that all isotropy groups are connected.  
It is known that when the symmetry rank \(k\) is maximal or almost maximal, \(M\) is diffeomorphic to a compact rank-one symmetric space.  
Investigate whether, under these intermediate symmetry assumptions, the following weaker conclusion must hold:  

> Is \(M\) *rationally elliptic*?  
>  
> More precisely, is it true that the total dimension of the rational homotopy groups of \(M\) is finite, i.e.  
> \[
> \sum_i \dim_\mathbb{Q}(\pi_i(M)\otimes \mathbb{Q}) < \infty,
> \]
> whenever \(M\) satisfies the geometric and symmetry hypotheses above?

Additionally, if rational ellipticity does hold, can one show further structural restrictions on the rational cohomology algebra of \(M\), analogous to those of torus manifolds or equal-rank homogeneous spaces?  

Why is it a "good" problem:  
This formulation isolates a central and tractable front in the intersection of **Riemannian geometry**, **transformation group theory**, and **rational homotopy theory**. It refines the broad conjectural landscape between the **Bott Conjecture** (“nonnegative curvature implies rational ellipticity”) and the **symmetry rank rigidity program** (which constrains the topology of highly symmetric manifolds).  

The problem seeks to determine whether **symmetry of intermediate rank** imposes the rational ellipticity property already conjectured in full generality. The threshold \(k \ge n/2\) mirrors known rigidity barriers in the Grove–Searle and Wilking theories, making the hypothesis geometrically meaningful.  

Resolving this question — even in special cases — would require integrating tools from **Sullivan minimal models** (rational homotopy data), **fixed-point set geometry** (via the slice theorem and Wilking’s connectedness lemma), and **equivariant cohomology** (used in torus manifold theory). This cross-disciplinary synthesis exemplifies how curvature, symmetry, and homotopy theory may jointly constrain manifold topology.  

The statement is deceptively simple — merely a geometric condition and a topological conclusion — yet it penetrates deeply into the structure of nonnegatively curved manifolds. A resolution would either provide new evidence for Bott’s conjecture under symmetry or produce counterexamples reshaping current beliefs, making it a "good" problem in the strongest sense.